\section{Topologies and completions}
\label{am-ch10-topology}
\source{atiyah-macdonald-commutative-algebra:page-0109}
\begin{definition}[Topological abelian group]
\label{am-def-topological-group}
\source{atiyah-macdonald-commutative-algebra:page-0110}
A topological abelian group has continuous addition and inversion.  Translation
is a homeomorphism, so the topology is determined by neighborhoods of $0$.
\end{definition}
\begin{lemma}[Hausdorff quotient]
\label{am-lem-10-1}
\dcref{10.1}
\source{atiyah-macdonald-commutative-algebra:page-0110}
If $H$ is the intersection of all neighborhoods of $0$ in a topological abelian
group $G$, then $H$ is a subgroup and the closure of $\{0\}$; $G/H$ is
Hausdorff; and $G$ is Hausdorff iff $H=0$.
\end{lemma}
\begin{proof}
Continuity of addition and inversion gives subgroup closure.  The condition
$x\in H$ is equivalent to $0\in x-U$ for every neighborhood $U$, which is
exactly $x\in\overline{\{0\}}$.  Cosets of $H$ are closed, so points of the
quotient are closed and the quotient is Hausdorff; the final equivalence follows.
\end{proof}
\begin{definition}[Completion and inverse limit]
\label{am-def-completion}
\source{atiyah-macdonald-commutative-algebra:page-0111}
For a countable neighborhood basis, the completion is the group of Cauchy
sequences modulo sequences tending to zero.  For a subgroup basis
$G=G_0\supseteq G_1\supseteq\cdots$, it is canonically
$\varprojlim G/G_n$, the coherent sequences under the quotient maps.  The
canonical map has kernel $\bigcap_nG_n$ and every continuous homomorphism
extends uniquely to completions.
\end{definition}
\begin{definition}[Inverse systems and subgroup topologies]
\label{am-def-inverse-system}
\source{atiyah-macdonald-commutative-algebra:page-0112}
For a decreasing subgroup basis $G=G_0\supseteq G_1\supseteq\cdots$, each
$G_n$ is both open and closed.  A Cauchy sequence determines an ultimately
constant class in every quotient $G/G_n$, and equivalent sequences determine
the same coherent family.  Conversely, a coherent family of classes has a
Cauchy representative.  More generally, an inverse system is a sequence of
groups $A_n$ with homomorphisms $\alpha_{n+1}:A_{n+1}\to A_n$; its inverse
limit consists of coherent sequences $(a_n)$ satisfying
$\alpha_{n+1}(a_{n+1})=a_n$.  An exact sequence of inverse systems is a
levelwise commutative diagram of exact sequences, and a system is surjective
when all transition maps are surjective.
\end{definition}
\begin{proposition}[Inverse limits of exact systems]
\label{am-prop-10-2}
\dcref{10.2}
\source{atiyah-macdonald-commutative-algebra:page-0113}
An exact sequence of inverse systems gives an exact sequence of inverse limits
at the first three terms.  If the left system has surjective transition maps,
the last map of inverse limits is also surjective.
\end{proposition}
\begin{proof}
For a system $A_n$ set $A=\prod_nA_n$ and
$d_A((a_n))=(a_n-\alpha_{n+1}(a_{n+1}))$.  Then
$\Ker d_A=\varprojlim A_n$.  Apply the snake lemma to the diagram of the three
products and their $d$ maps.  Surjectivity of the left transitions lets one
solve $x_n-\alpha_{n+1}(x_{n+1})=a_n$ inductively, so $d_A$ is surjective and
the connecting cokernel vanishes.
\uses{am-prop-2-10}
\end{proof}
\begin{remark}[Derived inverse-limit obstruction]
\label{am-rem-derived-limit}
\source{atiyah-macdonald-commutative-algebra:page-0113}
For the product differential $d_A((a_n))=a_n-\alpha_{n+1}(a_{n+1})$ used in
the proof of Proposition~\ref{am-prop-10-2}, the group $\Coker d_A$ is commonly
written $\varprojlim{}^{1} A_n$; it measures the failure of exactness at the
last inverse-limit term.
\uses{am-prop-10-2}
\end{remark}
\begin{corollary}[Completion of an exact sequence of groups]
\label{am-cor-10-3}
\dcref{10.3}
\source{atiyah-macdonald-commutative-algebra:page-0113}
For an exact sequence of groups with subgroup filtrations, the induced sequence
of completions is exact.
\uses{am-prop-10-2}
\end{corollary}
\begin{proof}
Apply Proposition~\ref{am-prop-10-2} to the exact quotient sequences
$0\to G'\cap G_n\to G_n\to pG_n\to0$.
\end{proof}
\begin{corollary}[Completion quotients]
\label{am-cor-10-4}
\dcref{10.4}
\source{atiyah-macdonald-commutative-algebra:page-0114}
The completion $\widehat G$ contains $G$ modulo its Hausdorff kernel and
$\widehat G/\widehat G_n\simeq G/G_n$.
\uses{am-cor-10-3,am-lem-10-1}
\end{corollary}
\begin{proof}
Apply Corollary~\ref{am-cor-10-3} to $G\to\widehat G$ and take inverse limits.
\end{proof}
\begin{proposition}[Completeness of the completion]
\label{am-prop-10-5}
\dcref{10.5}
\source{atiyah-macdonald-commutative-algebra:page-0114}
The completion $\widehat G$ is complete and Hausdorff.
\uses{am-cor-10-4,am-lem-10-1}
\end{proposition}
\begin{proof}
Taking inverse limits in Corollary~\ref{am-cor-10-4} identifies the completion
of $\widehat G$ with $\widehat G$ itself; Hausdorffness follows from
Lemma~\ref{am-lem-10-1}.
\end{proof}
\begin{definition}[Adic topology and completion]
\label{am-def-adic-completion}
\source{atiyah-macdonald-commutative-algebra:page-0114}
For an ideal $\aideal\subseteq A$, the $\aideal$-adic topology on $A$ has
$\aideal^n$ as neighborhoods of zero; on an $A$-module $M$ it has
$\aideal^nM$.  The completion is denoted $\widehat A$ or $\widehat M$.
The topology on $A$ is Hausdorff iff $\bigcap_n\aideal^n=0$.
\end{definition}

\section{Filtrations and graded rings}
\label{am-ch10-filtrations}
\begin{definition}[Filtration]
\label{am-def-filtration}
\source{atiyah-macdonald-commutative-algebra:page-0114}
An $\aideal$-filtration is $M=M_0\supseteq M_1\supseteq\cdots$ with
$\aideal M_n\subseteq M_{n+1}$; it is stable if equality holds for all large
$n$.  The powers $\aideal^nM$ form a stable filtration.
\end{definition}
\begin{lemma}[Bounded difference of stable filtrations]
\label{am-lem-10-6}
\dcref{10.6}
\source{atiyah-macdonald-commutative-algebra:page-0115}
Two stable $\aideal$-filtrations of a finitely generated module have bounded
difference: for some $n_0$, $M_{n+n_0}\subseteq M'_n$ and
$M'_{n+n_0}\subseteq M_n$ for every $n$.  They define the same topology.
\end{lemma}
\begin{proof}
It suffices to compare with $\aideal^nM$.  The filtration condition gives
$\aideal^nM\subseteq M_n$, while stability of $M'$ gives
$M'_{n+n_0}=\aideal^nM'_{n_0}\subseteq\aideal^nM$; interchange the two
filtrations for the other inclusion.
\end{proof}
\begin{definition}[Graded ring and module]
\label{am-def-graded}
\source{atiyah-macdonald-commutative-algebra:page-0115}
A graded ring is $A=\bigoplus_{n\ge0}A_n$ with $A_mA_n\subseteq A_{m+n}$;
a graded module is $M=\bigoplus M_n$ with $A_mM_n\subseteq M_{m+n}$.
Homogeneous components are unique.  The positive part $A_+=\bigoplus_{n>0}A_n$
is an ideal.  For an ideal filtration, $\operatorname{gr}_{\aideal}(A)=\bigoplus
\aideal^n/\aideal^{n+1}$ and similarly
$\operatorname{gr}(M)=\bigoplus M_n/M_{n+1}$.
\end{definition}
\begin{proposition}[Noetherian graded rings]
\label{am-prop-10-7}
\dcref{10.7}
\source{atiyah-macdonald-commutative-algebra:page-0115}
A graded ring $A$ is Noetherian iff $A_0$ is Noetherian and $A$ is finite type
as an $A_0$-algebra.
\end{proposition}
\begin{proof}
If $A$ is Noetherian, $A_0$ is a quotient and $A_+$ has finitely many
homogeneous generators; induction on degree shows they generate $A$.  The
converse is Hilbert's basis theorem.
\uses{am-thm-7-5}
\end{proof}
\begin{lemma}[Stable filtration criterion]
\label{am-lem-10-8}
\dcref{10.8}
\source{atiyah-macdonald-commutative-algebra:page-0116}
For Noetherian $A$, ideal $\aideal$, finite $M$, an $\aideal$-filtration is stable
iff its graded module is finite over $\operatorname{gr}_{\aideal}(A)$.
\end{lemma}
\begin{proof}
The graded pieces up to a fixed degree generate a graded submodule; these
submodules form an ascending chain.  Noetherianity makes the chain stop, which
is exactly eventual equality $M_{n+r}=\aideal^rM_n$.  Stability makes the
graded module generated by finitely many initial pieces.
\uses{am-prop-10-7,am-prop-6-2}
\end{proof}
\begin{proposition}[Artin--Rees lemma]
\label{am-prop-10-9}
\dcref{10.9}
\source{atiyah-macdonald-commutative-algebra:page-0116}
If $A$ is Noetherian, $M$ finite, $M'$ a submodule, and $(M_n)$ stable, then
$(M'\cap M_n)$ is a stable filtration of $M'$.
\end{proposition}
\begin{proof}
The intersections form an $\aideal$-filtration.  Their graded module embeds
in the Noetherian graded module of $M$, so Lemma~\ref{am-lem-10-8} gives
stability.
\uses{am-lem-10-8}
\end{proof}
\begin{corollary}[Artin--Rees equality]
\label{am-cor-10-10}
\dcref{10.10}
\source{atiyah-macdonald-commutative-algebra:page-0116}
For some $k$ and all $n\ge k$,
$M'\cap\aideal^nM=\aideal^{n-k}(M'\cap\aideal^kM)$.
\uses{am-prop-10-9}
\end{corollary}
\begin{proof}
Apply stability of the filtration $M'\cap\aideal^nM$.
\end{proof}
\begin{theorem}[Artin--Rees topological theorem]
\label{am-thm-10-11}
\dcref{10.11}
\source{atiyah-macdonald-commutative-algebra:page-0117}
For Noetherian $A$, finite $M$, and $M'\subseteq M$, the filtrations
$(\aideal^nM')$ and $(M'\cap\aideal^nM)$ have bounded difference; hence the
adic topology on $M'$ equals the topology induced from $M$.
\uses{am-prop-10-9,am-lem-10-6}
\end{theorem}
\begin{proof}
Artin--Rees makes the intersection filtration stable; Lemma~\ref{am-lem-10-6}
then compares it to the intrinsic power filtration.
\end{proof}

\section{Exactness and properties of completion}
\label{am-ch10-exact-completion}
\begin{proposition}[Exactness of finite completion]
\label{am-prop-10-12}
\dcref{10.12}
\source{atiyah-macdonald-commutative-algebra:page-0117}
The completion functor is exact on short exact sequences of finite modules over
a Noetherian ring.
\end{proposition}
\begin{proof}
Apply Corollary~\ref{am-cor-10-3} to the adic quotient sequences and use the
topological identification in Theorem~\ref{am-thm-10-11}.
\uses{am-cor-10-3,am-thm-10-11}
\end{proof}
\begin{proposition}[Completion as tensor product]
\label{am-prop-10-13}
\dcref{10.13}
\source{atiyah-macdonald-commutative-algebra:page-0117}
For finite $M$, the natural map $\widehat A\otimes_A M\to\widehat M$ is
surjective for every $A$; it is an isomorphism when $A$ is Noetherian.
\end{proposition}
\begin{proof}
Use a finite presentation $0\to N\to A^r\to M\to0$.  Completion commutes
with finite free sums, tensoring is right exact, and completion is exact for
finite modules by Proposition~\ref{am-prop-10-12}; the resulting diagram and
the five-lemma argument give surjectivity and then injectivity.
\uses{am-prop-10-12,am-prop-2-18,am-prop-2-3}
\end{proof}
\begin{proposition}[Flatness of the completed ring]
\label{am-prop-10-14}
\dcref{10.14}
\source{atiyah-macdonald-commutative-algebra:page-0118}
For Noetherian $A$, $\widehat A$ is a flat $A$-algebra.
\uses{am-prop-10-12,am-prop-10-13,am-prop-2-19}
\end{proposition}
\begin{proof}
The tensor-completion comparison is an isomorphism on finite modules, and
completion is exact there; the flatness criterion of Proposition~\ref{am-prop-2-19}
applies.
\end{proof}
\begin{proposition}[Powers and Jacobson radical after completion]
\label{am-prop-10-15}
\dcref{10.15}
\source{atiyah-macdonald-commutative-algebra:page-0118}
For Noetherian $A$ and its $\aideal$-adic completion,
$\widehat{\aideal}=\widehat A\otimes_A\aideal$,
$(\aideal^n)^{\widehat{}}=\widehat{\aideal}^{,n}$,
$\widehat A/\widehat{\aideal}^{,n}\simeq A/\aideal^n$, and
$\widehat{\aideal}\subseteq J(\widehat A)$.
\uses{am-prop-10-13,am-prop-10-5,am-prop-1-9}
\end{proposition}
\begin{proof}
Apply Proposition~\ref{am-prop-10-13} to the finite ideals $\aideal^n$ and
use inverse-limit quotients.  The geometric series
$(1-x)^{-1}=1+x+x^2+\cdots$ converges for $x\in\widehat{\aideal}$, so the
unit criterion Proposition~\ref{am-prop-1-9} gives the Jacobson assertion.
\end{proof}
\begin{proposition}[Completion of a local ring]
\label{am-prop-10-16}
\dcref{10.16}
\source{atiyah-macdonald-commutative-algebra:page-0118}
The completion of a Noetherian local ring $(A,\m)$ is local with maximal ideal
$\widehat\m$.
\uses{am-prop-10-15}
\end{proposition}
\begin{proof}
The residue quotient is unchanged and is a field; Proposition~\ref{am-prop-10-15}
places $\widehat\m$ in the Jacobson radical, so it is the unique maximal ideal.
\end{proof}
\begin{theorem}[Krull intersection theorem]
\label{am-thm-10-17}
\dcref{10.17}
\source{atiyah-macdonald-commutative-algebra:page-0119}
For Noetherian $A$, ideal $\aideal$, finite module $M$, the kernel of
$M\to\widehat M$ is the set of elements annihilated by some $1-a$ with
$a\in\aideal$.
\end{theorem}
\begin{proof}
Let $E=\bigcap_n\aideal^nM$.  The induced topology on $E$ is trivial, while
Artin--Rees identifies it with its $\aideal$-topology; hence $\aideal E=E$.
Nakayama's determinant corollary gives $(1-a)E=0$ for some $a\in\aideal$.
Conversely, $(1-a)x=0$ implies $x=a^nx\in\aideal^nM$ for all $n$.
\uses{am-thm-10-11,am-cor-2-5}
\end{proof}
\begin{corollary}[Krull intersection in a domain]
\label{am-cor-10-18}
\dcref{10.18}
\source{atiyah-macdonald-commutative-algebra:page-0119}
For a Noetherian domain and proper ideal $\aideal$, $\bigcap_n\aideal^n=0$.
\uses{am-thm-10-17}
\end{corollary}
\begin{proof}
In a domain $1+\aideal$ has no zero-divisors, so Theorem~\ref{am-thm-10-17}
gives zero kernel.
\end{proof}
\begin{corollary}[Hausdorff adic topology]
\label{am-cor-10-19}
\dcref{10.19}
\source{atiyah-macdonald-commutative-algebra:page-0119}
If $\aideal\subseteq J(A)$, the $\aideal$-topology on every finite module is
Hausdorff.
\uses{am-thm-10-17,am-prop-1-9}
\end{corollary}
\begin{proof}
Every $1-a$ is a unit by Proposition~\ref{am-prop-1-9}, so
Theorem~\ref{am-thm-10-17} gives zero intersection.
\end{proof}
\begin{corollary}[Maximal-ideal topology]
\label{am-cor-10-20}
\dcref{10.20}
\source{atiyah-macdonald-commutative-algebra:page-0120}
For a Noetherian local ring $(A,\m)$ and finite $M$, the $\m$-adic topology is
Hausdorff; in particular $\bigcap_n\m^n=0$ in $A$.
\uses{am-cor-10-19}
\end{corollary}
\begin{proof}
Take $\aideal=\m$ in Corollary~\ref{am-cor-10-19}.
\end{proof}
\begin{corollary}[Intersection of primary ideals]
\label{am-cor-10-21}
\dcref{10.21}
\source{atiyah-macdonald-commutative-algebra:page-0120}
For a Noetherian ring and prime $\p$, the intersection of all $\p$-primary
ideals is the kernel of $A\to A_\p$.
\uses{am-cor-10-20,am-prop-4-8}
\end{corollary}
\begin{proof}
Apply Corollary~\ref{am-cor-10-20} in $A_\p$ and lift through the primary
correspondence of Proposition~\ref{am-prop-4-8}.
\end{proof}

\section{Associated graded rings and completed Noetherian rings}
\label{am-ch10-graded-completion}
\begin{proposition}[Associated graded ring]
\label{am-prop-10-22}
\dcref{10.22}
\source{atiyah-macdonald-commutative-algebra:page-0120}
For Noetherian $A$ and ideal $\aideal$, $\operatorname{gr}_{\aideal}(A)$ is
Noetherian; it is canonically isomorphic to the associated graded ring of
$\widehat A$; and for a finite module with stable filtration,
$\operatorname{gr}(M)$ is finite over it.
\end{proposition}
\begin{proof}
If $x_1,\ldots,x_r$ generate $\aideal$, the graded ring is generated over
$A/\aideal$ by their degree-one classes, so use Hilbert's theorem.  The graded
piece isomorphisms follow from Proposition~\ref{am-prop-10-15}; finite
generation follows from stability and finite generation of the initial pieces.
\uses{am-thm-7-5,am-prop-10-15,am-lem-10-8}
\end{proof}
\begin{lemma}[Filtered maps]
\label{am-lem-10-23}
\dcref{10.23}
\source{atiyah-macdonald-commutative-algebra:page-0121}
For a homomorphism of filtered groups, injectivity (respectively surjectivity)
of the associated graded map implies injectivity (respectively surjectivity) of
the map on completions.
\end{lemma}
\begin{proof}
Apply the snake lemma to
$0\to A_n/A_{n+1}\to A/A_{n+1}\to A/A_n\to0$ and induct on $n$; the
surjective case also gives surjective transition maps on kernels, so
Proposition~\ref{am-prop-10-2} applies.
\uses{am-prop-2-10,am-prop-10-2}
\end{proof}
\begin{proposition}[Finite generation from the graded module]
\label{am-prop-10-24}
\dcref{10.24}
\source{atiyah-macdonald-commutative-algebra:page-0121}
If $A$ is complete in its $\aideal$-topology, $M$ is Hausdorff in a filtration,
and $\operatorname{gr}(M)$ is finite over $\operatorname{gr}(A)$, then $M$ is
finite over $A$.
\end{proposition}
\begin{proof}
Lift finitely many homogeneous graded generators to elements of $M$.  Map a
finite direct sum of shifted copies of $A$ to $M$.  The graded map is surjective,
so Lemma~\ref{am-lem-10-23} makes the completion map surjective.  Completeness
of $A$ and Hausdorffness of $M$ identify the completed map with the original,
so the lifts generate.
\uses{am-lem-10-23}
\end{proof}
\begin{corollary}[Noetherianity from the graded module]
\label{am-cor-10-25}
\dcref{10.25}
\source{atiyah-macdonald-commutative-algebra:page-0122}
Under the hypotheses of Proposition~\ref{am-prop-10-24}, if
$\operatorname{gr}(M)$ is Noetherian then $M$ is Noetherian.
\uses{am-prop-10-24,am-prop-6-2}
\end{corollary}
\begin{proof}
For $M'\subseteq M$, the induced graded module embeds in $\operatorname{gr}(M)$
and is finite; it is Hausdorff, so Proposition~\ref{am-prop-10-24} makes $M'$ finite.
\end{proof}
\begin{theorem}[Noetherianity of completion]
\label{am-thm-10-26}
\dcref{10.26}
\source{atiyah-macdonald-commutative-algebra:page-0122}
The adic completion of a Noetherian ring is Noetherian.
\end{theorem}
\begin{proof}
Its associated graded ring is Noetherian by Proposition~\ref{am-prop-10-22};
apply Corollary~\ref{am-cor-10-25} to the complete ring filtered by powers of
the completed ideal.
\uses{am-prop-10-22,am-cor-10-25}
\end{proof}
\begin{corollary}[Formal power series]
\label{am-cor-10-27}
\dcref{10.27}
\source{atiyah-macdonald-commutative-algebra:page-0122}
If $A$ is Noetherian, then $A[[x_1,\ldots,x_n]]$ is Noetherian.
\uses{am-thm-10-26,am-cor-7-6}
\end{corollary}
\begin{proof}
It is the completion of $A[x_1,\ldots,x_n]$ in the ideal generated by the
variables; use Theorem~\ref{am-thm-10-26}.
\end{proof}
